\section{Resources and indicative timeline}
\label{sec:resources}

\begin{evfutbox}
  Every figure below is a planning range, not an estimate fitted to data: nothing in the
  validation program (\cref{sec:validation}) has run. Personnel cost per person-month (PM) varies
  by a factor of two or more across European countries and institution types. The range used
  here, EUR~6{,}000--10{,}000 per PM fully loaded, should be replaced with the host institution's
  own rate once a host is named (\cref{sec:team}).
\end{evfutbox}

\subsection{Money released by each gate}

The budget is split into tranches, and each tranche is released only when the gate before it
reads continue (\cref{sec:roadmap}). \Cref{tab:resources-totals} gives each tranche's gate, scope,
effort and cost. Phase~2 has two lines. The main line (2a--2c) carries H2, H3, the partner
pilot and the learner link; the organizational line (O) carries H-OSS. Phase~0 is owner-run and
needs no grant.

\begin{table}[htbp]
  \centering
  \footnotesize
  \caption{Tranches, the gate that releases each, and indicative cost. PM at EUR 6--10k; other
    direct costs are participants, coders, compute, legal, open access, travel and the gate
    reader's fee. Months count from the start of Phase~1.}
  \label{tab:resources-totals}
  \begin{tabularx}{\linewidth}{@{}>{\raggedright\arraybackslash}p{0.1\linewidth}>{\raggedright\arraybackslash}p{0.11\linewidth}>{\raggedright\arraybackslash}X>{\raggedright\arraybackslash}p{0.1\linewidth}p{0.07\linewidth}p{0.09\linewidth}p{0.11\linewidth}@{}}
    \toprule
    Tranche & Released by & Scope & Months & PM & Other direct & Total (EUR) \\
    \midrule
    Phase~0 & Owner, before any grant & V1 closure study (WP3) & 2--4 mo before M1 & Owner time &
      $\approx10^3$ & $\approx10^3$ \\
    Phase~1 & MS1 continue & Area score, freeze, V0, judge text-type check, pooled V2, V4 (WP1--4,
      WP11) & M1--M14 & 35--44 & 20k--70k & \textbf{0.23--0.51\,M} \\
    2a & MS3 continue & Workbench, V5 pilot, V8a, ledger groundwork & M15--M24 & 29--44 &
      15k--50k & 0.19--0.49\,M \\
    2b & MS4 continue & V5 main, powered; high end at the 450-area ceiling & M25--M38 & 22--36 &
      35k--250k & 0.17--0.61\,M \\
    2c & MS5 continue & V6, V7, V8c, education pilot, on-prem stack & M39--M48 & 38--62 &
      35k--165k & 0.26--0.79\,M \\
    O & MS3 continue & Git adapter, code lenses, V3 (H-OSS) & M15--M36 & 14--22 & 5k--35k &
      0.09--0.26\,M \\
    \midrule
    \textbf{Phase~2} & & 2a--2c and O & M15--M48 & 103--164 & 90k--500k &
      \textbf{0.7--2.2\,M} \\
    \textbf{Program} & & Phases~1 and 2 & 48~mo & 138--208 & 110k--570k & \textbf{0.9--2.7\,M} \\
    \midrule
    Interview-only pivot & MS1, MS3 or MS5 stop & Separate line (see below) & 18--24~mo & 21--31 &
      15k--70k & 0.14--0.38\,M \\
    Product (S6) & MS5 and MS7 & Team of 8--15 & 18--36~mo later & -- & Hosting, sales &
      Outside this estimate \\
    \bottomrule
  \end{tabularx}
\end{table}

\evint{} At EUR~8k per PM, Phase~1 costs EUR~0.30--0.42\,M; the full range is EUR~0.23--0.51\,M,
so about EUR~0.3--0.5\,M. Every Phase~2 euro waits for MS3.

\paragraph{Money at risk if H2 fails.} The spend on H2 before a failure is visible depends on the
gate that reads it.
\begin{itemize}[noitemsep]
  \item \textbf{At MS1:} about EUR~$10^3$ of coder time and compute, plus the owner's time.
  \item \textbf{At MS3:} Phase~1, EUR~0.23--0.51\,M. It still leaves a released closure set, a
    frozen area score and a public reconstruct-versus-CTA benchmark.
  \item \textbf{At MS4:} Phase~1 plus tranche 2a, EUR~0.42--1.00\,M.
  \item \textbf{At MS5:} Phase~1 plus 2a and 2b, EUR~0.59--1.61\,M. Tranche 2c, EUR~0.26--0.79\,M,
    is never released.
\end{itemize}
The organizational line is outside these sums because H-OSS does not depend on H2. It runs to its
own gate whatever H2's gates read, and it releases no partner pilot on its own.

\paragraph{The interview-only pivot, costed separately.} If MS1 stops the automated absence check,
or MS3 or MS5 ends H2 funding, the program can propose one smaller study. It tests targeted
interviewing guided by expected information gain against untargeted interviewing, with no
reconstruction prior. It needs a light workbench (8--12 PM), the efficiency study in
interview-only form (10--14 PM), ethics and pre-registration (2--3 PM), and release (1--2 PM).
Participants and coders add EUR~15k--70k. The line is a new proposal read by the gate reader; it never draws on Phase~2's
budget.

\subsection{Person-months per work package}

\Cref{tab:resources-pm} gives effort by work package and phase. Coder hours are counted under
participant costs (\cref{tab:resources-participants}), not here, because they are hourly
contracts. Phase~1's 35--44 PM over 12--15 months fit a team of 2.5--3.0~FTE; Phase~2's 103--164
PM over 34 months average 3.0--4.8~FTE (\cref{sec:team}). V5's main study concentrates about
22--36 PM in 14 months, which is why the CTA practitioner rises to full time in S3.

\begin{table}[htbp]
  \centering
  \footnotesize
  \caption{Person-months per work package and phase.}
  \label{tab:resources-pm}
  \begin{tabularx}{\linewidth}{@{}Xp{0.1\linewidth}p{0.1\linewidth}p{0.3\linewidth}@{}}
    \toprule
    Work package & Phase~1 & Phase~2 & Main roles \\
    \midrule
    WP1 Pre-registration, governance, data protection & 4--5 & 6--10 &
      PI, project manager, statistician, counsel \\
    WP2 Area score, freeze, deferred N2 validation (V0) & 10--13 & -- &
      Research engineer \\
    WP3 Absence-check validation (V1 in P0; text-type check) & 4--6 & -- &
      CTA practitioner, statistician, engineer \\
    WP4 Pooled CTA gold test and residual check (V2, V4) & 15--18 & -- &
      Engineer, research assistant, statistician \\
    WP5 Organizational track: H-OSS (V3) & -- & 14--22 &
      AI/ML researcher, platform engineer, statistician \\
    WP6 Elicitation workbench & -- & 16--25 &
      Software engineer, frontend/HCI, CTA practitioner \\
    WP7 Prospective expert validation (V5, V6) & -- & 31--49 &
      CTA practitioner, interviewers, statistician \\
    WP8 Organizational pilot (V7) & -- & 10--16 &
      Project manager, security, counsel, CTA practitioner \\
    WP9 Learner link and education pilot & -- & 10--16 &
      Learning scientist, statistician \\
    WP10 Platform groundwork & -- & 12--20 &
      Platform engineer, security engineer \\
    WP11 Open benchmark and dissemination & 2 & 4--6 &
      PI, research engineer \\
    \midrule
    \textbf{Total} & \textbf{35--44} & \textbf{103--164} & \\
    \bottomrule
  \end{tabularx}
\end{table}

\subsection{Compute, API budgets, and the headroom lesson}

\evobs{} The PoC gives one hard number to plan from: a single reconstruction run cost \$0.68 (PLC)
or \$0.50 (GDPR).\footnote{\repo{research/n2/e_plant_eabst_report.md} \S1.2, \S3;
\repo{research/n3/README.md}.} Money was not the constraint; \textbf{headroom} was. A weekly key
limit of \$5, set below the sum of the experiments' registered caps (about \$21), blocked E-LIVE~v3 at its first
call.\footnote{\repo{docs/lessons.md}.} The budget stop is one reason N2's gated arms were never
scored (\cref{sec:n2}).

\evint{} The lesson has four parts. (1)~Provider spending limits are set at least 3$\times$ the sum
of a stage's registered caps, on a monthly window, and recorded in the protocol before the first
run. (2)~The headroom check runs as code before a run sequence starts. (3)~Each experiment has its
own API key, so one study cannot exhaust another's budget. (4)~Every study reports cost per
\emph{validated} item, not cost per claim.

Later stages cost more per run than the PoC did. A decisive test needs at least three reruns per
task, because retrieval variance dominated the PoC's cross-run instability. It also needs two or
three model families for the generator, verifier and closure judge, stronger models, plain-LLM
baselines, memorization probes and, on the organizational line, repository histories.
Compute stays small next to personnel: about EUR~$10^3$--$10^4$ in Phase~1 and
EUR~$10^3$--$10^4$ per tranche in Phase~2. \Cref{tab:resources-compute} in \cref{app:program}
gives the orders of magnitude by phase.


\subsection{Participant, coder and expert-hour costs}

\Cref{tab:resources-participants} in \cref{app:program} gives honoraria or wages, not
consulting rates, at the volumes set by \cref{sec:validation}. Expert time is costed at EUR~60--200 per hour and coder time at
EUR~20--45 per hour. Over Phase~2, participant, coder and learner costs run roughly EUR~55k--355k,
the high end at V5's funded ceiling. Other direct costs (legal counsel, compute, open access,
travel, ethics fees, the gate reader) add roughly EUR~35k--140k.

\paragraph{V5 is budgeted for power.} \evint{} Deciding H2 at $\delta = 0.05$ needs about 300
areas, about 150 per class (\cref{sec:validation}). With 3 experts per area, that is about 900
probes of about 8 minutes each. At about 20 areas per expert, V5 needs about 45 experts, about 22
per domain, plus 6 in the pilot. Each expert gives about 3 hours in sessions and 1 hour rating
other experts' items: about 150--250 expert-hours in all, and about 70 session hours per
interviewer. Coding costs more than the experts: about 120 hours of audio need 1,200--1,900 coder
hours. Honoraria, recruitment and coding come to roughly EUR~40k--160k before staff time. MS4
releases the main study only if the pilot's projection fits a funded ceiling of 450 areas, so
tranche 2b's high end is costed at that ceiling. At a prevalence of 0.3, about 500 areas would be
needed; that exceeds the ceiling, and MS4 would not release V5. The binding constraint is
recruiting about 45 eligible experts, not money.


\subsection{Timeline uncertainty}

\evint{} Four facts drive most of the uncertainty, and none can be resolved before the program
starts:
\begin{itemize}[noitemsep]
  \item the number of itemized CTA golds that can be obtained and dated (fewer golds mean fewer
    areas, which makes CONTINUE at MS3 harder to reach);
  \item the prevalence of residual-present areas and the coding time per probe in V5 (they
    scale areas, expert-hours and coder hours);
  \item the number of usable developer departures for H-OSS;
  \item the host institution's rate, and whether V7's legal work takes weeks or months.
\end{itemize}
A stop at MS1, MS3 or MS5 ends the spending that depends on it, and the stop is reported as a
completed result.

\subsection{Funding instrument fit, by type}

No specific call or deadline is named. Eligibility depends on a host institution, which the
program does not yet have. \evfut{} \Cref{tab:resources-funding}, in \cref{app:program}, matches
phases to instrument types only.

\evint{} The strongest fit for Phase~1 is a single-institution methods grant, submitted only after
MS1 reads continue, with the V1 numbers on its first page. It buys the pooled screening result
(MS3), the frozen configuration and area score, the judge's text-type check, and a public
benchmark. It does not buy a verdict on H2; only MS5 gives one.
